\bookchapter{Chapitre 15 : Science de la batterie et optimisation de l'autonomie}

\chapterintro{La batterie est le « cœur énergétique » de votre smartphone. Ce chapitre vous aide à comprendre les petites humeurs de la batterie, à apprendre des techniques scientifiques de charge et d'autonomie, et à dire adieu à l'anxiété de la batterie.}

\booksection{I. Comment fonctionne la batterie ? Cycles et état de santé}

La batterie d'un smartphone (lithium-ion ou lithium-polymère) est comme une éponge d'énergie : la charger c'est absorber de l'eau (stocker l'électricité), la décharger c'est presser l'eau (utiliser l'électricité). Quelques caractéristiques à retenir :

\begin{itemize}
  \item \textbf{Charge à la demande}: contrairement aux anciennes batteries, elle n'a pas d'effet mémoire, pas besoin de la vider complètement avant de recharger.
  \item \textbf{Peur de la chaleur}: les températures élevées sont le tueur numéro un des batteries.
  \item \textbf{Vieillissement}: avec le temps, la capacité maximale diminue, c'est un phénomène normal.
\end{itemize}

\booksubsection{Indicateurs clés : nombre de cycles et état de santé}

\begin{itemize}
  \item Un cycle n'est pas une charge. Par exemple, si aujourd'hui vous utilisez 60 \% puis chargez à plein, demain 40 \% puis chargez à plein, la consommation totale de 100 \% compte pour un cycle.
  \item L'état de santé de la batterie est le pourcentage de la capacité maximale actuelle par rapport à celle neuve. En dessous de 80 \%, vous sentirez clairement l'autonomie diminuer.
  \item \textbf{Comment vérifier ?}iPhone : Réglages > Batterie > Santé de la batterie et charge ; Android : généralement recherchez « santé de la batterie » dans les réglages.
\end{itemize}

\booksection{II. Mythes de la charge ? Charge rapide, charge sans fil et règle d'or}

\begin{itemize}
  \item \textbf{La charge rapide abîme-t-elle la batterie ?}Avec des accessoires de charge rapide d'origine ou certifiés, les dommages sont minimes. Le téléphone et le chargeur ont un contrôle de température, ils réduisent automatiquement la puissance en cas de surchauffe.
  \item \textbf{Charge sans fil}: pratique mais moins efficace, chauffe davantage. Adaptée pour la charge de toute la nuit.
  \item \textbf{Règle d'or de préservation}: essayez de maintenir le niveau de charge entre 20 \% et 80 \%, ne la laissez pas affamée (0 \%) ou gavée (100 \%).
\end{itemize}

\begin{tcbinfo}{Fonction « limite de charge 80 \% » de l'iPhone — transformer la règle d'or en paramètre système}
Les séries iPhone 15/16 (iOS 17 et supérieur) ont intégré la « règle des 20 \%-80 \% » en tant que fonctionnalité système. Une fois activée, le smartphone cesse automatiquement de charger à 80 \%, retardant considérablement le vieillissement de la batterie. Chemin de paramétrage :

\texttt{Réglages > Batterie > Santé de la batterie et charge > Optimisation de la charge > Limite 80 \%}

Signification des trois options : (1)\textbf{Sans limite}— charge jusqu'à 100 \% ; (2)\textbf{Optimisation de la charge}— apprend vos habitudes et ne retarde la charge à 100 \% que pendant la dernière partie de la nuit ; (3)\textbf{Limite 80 \%}— s'arrête directement à 80 \%, le plus protecteur pour la batterie. Pour ceux qui ont constamment accès à une prise à domicile ou au bureau, la limite 80 \% est fortement recommandée ; pour ceux qui voyagent souvent et ont besoin d'une batterie pleine, l'optimisation de la charge est préférable.

\end{tcbinfo}

\begin{tcbinfo}{Chargeur GaN — le meilleur allié des voyageurs}
Vous avez peut-être remarqué que les chargeurs rapides sont de plus en plus petits, grâce à\textbf{GaN (nitrure de gallium)}ce nouveau matériau semi-conducteur. Par rapport aux chargeurs traditionnels à base de silicium, les chargeurs GaN présentent plusieurs avantages :

\begin{itemize}
  \item \textbf{Taille réduite}: À puissance égale, un chargeur GaN peut faire 1/2 voire 1/3 de la taille d'un chargeur silicium, un chargeur GaN 65 W est plus petit que l'ancien chargeur 5 W d'Apple.
  \item \textbf{Rendement élevé, faible échauffement}: Les pertes de commutation du nitrure de gallium sont bien inférieures à celles du silicium, le rendement de conversion énergétique peut dépasser 95 \%, le chargeur ne chauffe pas facilement.
  \item \textbf{Idéal pour les voyages}: Un chargeur GaN multiport 65 W ou 100 W peut charger simultanément un smartphone, une tablette et un ordinateur portable, beaucoup plus léger que de transporter trois « briques ».
  \item \textbf{Conseil d'achat}: Choisissez un chargeur GaN prenant en charge le protocole PD (65 W minimum pour un ordinateur portable), des marques connues comme Baseus, Anker, Nubia, Xiaomi, etc. conviennent. Vérifiez que le câble de données est compatible (les hautes puissances nécessitent un câble USB-C avec puce E-Marker).
\end{itemize}

\end{tcbinfo}

\booksection{III. Analyse complète des protocoles de charge rapide : pourquoi votre chargeur ne charge-t-il pas vite ?}

Vous avez peut-être rencontré cette situation : votre smartphone prend en charge la charge rapide, mais après avoir changé de chargeur, la charge devient très lente. C'est parce que la charge rapide n'est pas simplement « plus de puissance suffit », mais nécessite que le chargeur, le câble et le téléphone « parlent la même langue » — ce sont les protocoles de charge rapide.

\booksubsection{Principe de base de la charge rapide}

Puissance de charge = Tension × Courant (P = U × I). Pour augmenter la puissance, on augmente soit la tension, soit le courant, soit les deux. Les différents protocoles ont choisi différentes approches technologiques.

Processus de la charge rapide :

\begin{enumerate}
  \item \textbf{Phase de handshake}: le téléphone et le chargeur communiquent via le câble de données, échangent leurs protocoles et informations de puissance pris en charge.
  \item \textbf{Phase de négociation}: les deux parties déterminent le niveau de puissance maximal communément pris en charge.
  \item \textbf{Phase de charge}: le chargeur produit la tension et le courant négociés, le téléphone gère l'alimentation côté récepteur.
  \item \textbf{Phase d'ajustement}: à mesure que le niveau de charge augmente, le téléphone réduit progressivement la puissance de charge pour protéger la batterie.
\end{enumerate}

\booksubsection{Détails des principaux protocoles de charge rapide}

\textbf{USB PD (Power Delivery)}— le roi des protocoles universels

\begin{itemize}
  \item \textbf{Caractéristiques}: standard officiel développé par l'USB-IF, la couverture la plus large. PD 3.0 prend en charge jusqu'à 100 W (20 V/5 A), PD 3.1 étendu à 240 W (48 V/5 A).
  \item \textbf{Paliers de tension}: 5 V, 9 V, 15 V, 20 V, PD 3.1 prend également en charge 28 V, 36 V, 48 V.
  \item \textbf{Compatibilité}: tous les iPhone, Google Pixel, certains smartphones Android haut de gamme, ordinateurs portables, la console Switch prennent en charge le protocole PD.
  \item \textbf{Avantages}: forte interopérabilité, un seul chargeur peut charger plusieurs appareils ; prend en charge PPS (Power Supply programmable), permet un réglage fin de la tension et du courant.
  \item \textbf{Inconvénients}: nécessite un port Type-C et une puce dédiée, coût plus élevé.
\end{itemize}

\textbf{Qualcomm QC (Quick Charge)}— le pionnier du camp Android

\begin{itemize}
  \item \textbf{QC 3.0}: prend en charge une tension réglable de 3,6 V à 20 V par paliers de 0,2 V, jusqu'à 18 W.
  \item \textbf{QC 4.0/4+}: compatible avec USB PD, prend en charge jusqu'à 27 W, intègre la technologie de charge double voie.
  \item \textbf{QC 5.0}: prend en charge plus de 100 W, 0-50 \% en 5 minutes seulement, améliore considérablement la gestion thermique.
  \item \textbf{Compatibilité}: la plupart des téléphones équipés de processeurs Qualcomm Snapdragon le prennent en charge, mais chaque marque privilégie généralement son propre protocole propriétaire.
\end{itemize}

\textbf{OPPO VOOC / OnePlus WARP / realme Dart}— le représentant de la basse tension et fort courant

\begin{itemize}
  \item \textbf{Approche technologique}: basse tension et fort courant (à partir de 5 V/4 A), moins de chauffe, meilleure efficacité.
  \item \textbf{VOOC 1.0}: 20 W (5 V/4 A), lancé en 2014, a inauguré l'ère de la charge rapide.
  \item \textbf{SuperVOOC}: 50 W-65 W, adopte une conception à double cellule en série.
  \item \textbf{SuperVOOC 2.0}: 65 W, charge une batterie de 4000 mAh en 30 minutes.
  \item \textbf{150 W/240 W Super Flash Charge}: la technologie de charge rapide de production la plus rapide actuellement.
  \item \textbf{Caractéristiques}: nécessite un chargeur et un câble dédiés (avec des contacts supplémentaires), incompatible avec d'autres protocoles.
\end{itemize}

\textbf{Huawei SuperCharge}— l'évolution de la charge rapide haute tension

\begin{itemize}
  \item \textbf{FCP (Fast Charge Protocol)}: 9 V/2 A, 18 W, compatible QC 2.0.
  \item \textbf{SCP (Super Charge Protocol)}: prend en charge une puissance plus élevée, de 22,5 W à 66 W puis 100 W.
  \item \textbf{Caractéristiques technologiques}: adopte la technologie de pompe de charge, convertit la haute tension en basse tension à fort courant pour la batterie, efficacité jusqu'à 98 \%.
  \item \textbf{Compatibilité}: réservé aux téléphones Huawei/Honor, certains chargeurs prennent en charge le protocole PD.
\end{itemize}

\textbf{Xiaomi Surge Charge rapide}— l'essor des protocoles propriétaires

\begin{itemize}
  \item \textbf{Puce Surge P1}: puce de gestion de charge rapide développée par Xiaomi, prend en charge la charge rapide filaire 120 W.
  \item \textbf{Puce Surge G1}: puce de gestion de batterie, combinée à la P1 pour une charge rapide plus efficace.
  \item \textbf{Caractéristiques technologiques}: adopte la solution double cellule en série + pompe de charge, prend en charge la compatibilité multi-protocoles.
\end{itemize}

\textbf{Charge rapide Apple}— conservatrice mais stable

\begin{itemize}
  \item \textbf{Protocole}: prend en charge uniquement le protocole USB PD.
  \item \textbf{Puissance}: iPhone 8 et versions ultérieures prennent en charge la charge rapide, jusqu'à environ 27 W (iPhone 15 Pro Max peut atteindre 27 W+).
  \item \textbf{Caractéristiques}: puissance relativement conservatrice, mais courbe de charge optimisée, excellente gestion thermique.
  \item \textbf{Exigences d'accessoires}: nécessite un câble USB-C vers Lightning (ou USB-C vers USB-C) et un chargeur PD.
\end{itemize}

\booksubsection{Tableau comparatif des protocoles de charge rapide}

\begin{center}
\begin{tabularx}{\linewidth}{|X|X|X|X|}
\hline
\textbf{Protocole} & \textbf{Puissance max} & \textbf{Approche technologique} & \textbf{Compatibilité} \tabularnewline
\hline
USB PD 3.1 & 240 W & Haute tension / multi-paliers & Le plus large \tabularnewline
QC 5.0 & 100 W+ & Haute tension réglable & Modèles Qualcomm \tabularnewline
VOOC/Dart & 240 W & Basse tension fort courant & Spécifique à OPPO et marques liées \tabularnewline
SuperCharge & 100 W & Pompe de charge & Huawei/Honor \tabularnewline
Charge rapide Apple & \textasciitilde{}27 W & Protocole PD & iPhone/iPad \tabularnewline
\hline
\end{tabularx}
\end{center}

\booksubsection{Protocoles de charge sans fil}

La charge sans fil a également ses propres standards de protocole :

\begin{itemize}
  \item \textbf{Standard Qi}: développé par le WPC (Wireless Power Consortium), le standard le plus universel. Le Qi de base prend en charge 5 W-15 W.
  \item \textbf{Standard Qi2}: publié en 2023, introduit le positionnement magnétique (similaire au MagSafe d'Apple), prend en charge 15 W+.
  \item \textbf{Protocoles propriétaires}: chaque marque a sa propre charge sans fil haute vitesse, comme la charge sans fil 50 W de Xiaomi, la charge sans fil 50 W de Huawei, etc., nécessitent un chargeur dédié.
\end{itemize}

\booksubsection{Guide d'achat de charge rapide}

\begin{itemize}
  \item \textbf{Vérifier la compatibilité des protocoles}: lors de l'achat d'un chargeur, confirmez qu'il prend en charge le protocole de charge rapide de votre téléphone. Les chargeurs tiers indiquent généralement les protocoles pris en charge.
  \item \textbf{Vérifier le câble de données}: la charge rapide à haute puissance nécessite un câble dédié. Les câbles ordinaires peuvent ne prendre en charge que moins de 60 W.
  \item \textbf{Vérifier la certification}: les utilisateurs Apple devraient choisir des câbles certifiés MFi ; les utilisateurs Android devraient choisir des câbles de marque d'origine ou de tiers reconnus.
  \item \textbf{Vérifier la chauffe}: il est normal que le téléphone chauffe pendant la charge rapide, mais s'il devient brûlant, ce n'est pas normal, il est conseillé de suspendre la charge ou de changer de chargeur.
\end{itemize}

\begin{tcbwarn}{Conseils de sécurité pour la charge rapide}
\begin{itemize}
  \item Évitez de jouer à des jeux gourmands pendant la charge, cela aggrave la chauffe.
  \item N'utilisez pas de chargeurs et de câbles de mauvaise qualité, cela peut présenter des risques de sécurité.
  \item Retirez les coques de téléphone épaisses pendant la charge pour aider à la dissipation thermique.
  \item Si le chargeur ou le téléphone chauffe anormalement ou dégage une odeur, arrêtez immédiatement de l'utiliser.
\end{itemize}

\end{tcbwarn}

\booksection{IV. Nouvelle technologie de batterie : anode en silicium-carbone}

Si vous avez suivi les présentations de smartphones phares en 2024-2026, le terme « anode silicium-carbone » est presque devenu synonyme de batterie haut de gamme. Il s'agit d'une mise à niveau importante du matériau d'anode des batteries lithium-ion, dont le principe n'est pas complexe :

\begin{itemize}
  \item \textbf{Anode en graphite traditionnelle}: Les batteries de smartphones ont toujours utilisé du graphite comme anode, technologie mature et stable, mais la densité énergétique a atteint sa limite théorique (environ 372 mAh/g).
  \item \textbf{Incorporation de silicium}: La capacité massique théorique du silicium est environ 10 fois supérieure à celle du graphite (environ 4200 mAh/g). En incorporant une petite quantité de silicium dans le graphite pour former une « anode composite silicium-carbone », la densité énergétique de la batterie peut être significativement augmentée.
  \item \textbf{Augmentation de la densité énergétique}: À volume égal, la capacité de la batterie à anode silicium-carbone peut augmenter d'environ\textbf{15\%-20\%}. C'est la raison clé pour laquelle les smartphones phares de 2024-2026 peuvent intégrer des batteries de 5500-6000 mAh dans des corps de 7-9 mm d'épaisseur.
\end{itemize}

\begin{tcbinfo}{Quels modèles l'utilisent ?}
\begin{itemize}
  \item \textbf{Honor}: La série Magic a été la première à commercialiser l'anode silicium-carbone, Magic6 Pro, Magic7, et le pliable Magic V3 l'utilisent tous.
  \item \textbf{Xiaomi}: Xiaomi 14 Ultra et la série Xiaomi 15 sont équipées de batteries de nouvelle génération à anode silicium-carbone.
  \item \textbf{vivo}: Les séries X100 Ultra et X200 Pro utilisent une solution à anode silicium-carbone à haute densité.
  \item \textbf{Huawei/OPPO, etc.}: Les modèles phares de 2025-2026 suivent généralement la tendance, l'anode silicium-carbone passe de « technologie de pointe » à équipement standard des modèles haut de gamme.
\end{itemize}

\end{tcbinfo}

\textbf{Impact pratique pour l'utilisateur}: Le plus grand avantage de l'anode silicium-carbone est qu'« à volume égal, la capacité de la batterie est plus importante » — le smartphone peut être plus fin sans sacrifier l'autonomie ; ou le corps garde la même épaisseur mais la batterie passe de 4500 mAh à 5500 mAh. Autrement dit, vous n'avez plus à payer le prix de « l'épaisseur et du poids » pour l'autonomie. À noter que les batteries à anode silicium-carbone nécessitent encore une optimisation avec la puce de charge rapide pour la gestion des basses températures et de la charge complète, c'est l'une des raisons pour lesquelles les fabricants recommandent de maintenir la charge à 80 \% à long terme.

\booksection{V. Les vraies astuces d'économie d'énergie}

\begin{itemize}
  \item \textbf{Écran (le roi de la consommation)}: baissez la luminosité, raccourcissez le délai de verrouillage automatique. Les téléphones à écran OLED économisent beaucoup d'énergie en mode sombre.
  \item \textbf{Signal faible}: dans les ascenseurs, les sous-sols, le téléphone cherche désespérément le signal, consommation très rapide.
  \item \textbf{5G}: si le signal est mauvais, passez manuellement en 4G, l'autonomie sera bien meilleure.
  \item \textbf{Services de localisation}: pour les applications non liées à la navigation, réglez la permission de localisation sur « autoriser uniquement pendant l'utilisation ».
  \item \textbf{Limiter l'activité en arrière-plan}: iPhone dans Réglages > Général > Actualisation en arrière-plan ; Android dans les paramètres de batterie des applications.
\end{itemize}

\begin{tcbpractice}{Faites un check-up : examinez votre batterie}
\begin{enumerate}
  \item Vérifiez l'état de santé : trouvez et notez la capacité maximale actuelle de votre batterie.
  \item Analysez les détails de consommation : regardez les graphiques et la liste des applications dans les réglages de batterie, identifiez les coupables.
  \item Élaborez une stratégie : limitez les applications qui consomment anormalement en arrière-plan.
\end{enumerate}

\end{tcbpractice}
